\section{Related Work}
\label{sec:related}

\paragraph{Partisan redistricting metrics.}
The efficiency gap \citep{stephanopoulos2015} measures partisan asymmetry as
the difference in wasted votes between the two parties.
Mean-median difference \citep{mcdonald2015} measures whether the median district
outcome differs systematically from the mean.
Both metrics characterise the \emph{outcome} of a redistricting plan but do
not prescribe an algorithm that achieves a target outcome.
\citet{mcghee2014} shows that geographic sorting produces Republican-favourable
gaps under most compactness criteria, even without intentional gerrymandering ---
the direct precursor to the Rodden effect documented empirically in B.9.

\paragraph{Proportional representation mechanisms.}
Multi-member districts and ranked-choice voting achieve proportionality by
changing the electoral rule rather than the map \citep{lijphart1994}.
Seat-optimal gerrymandering \citep{chen2013, friedman2009} uses partisan data
to draw maximally advantageous maps.
Our approach differs from both: we use partisan data only to the minimum extent
needed for proportionality, while retaining the single-member, first-past-the-post
structure and the minimum-edge-cut geographic objective.

\paragraph{Algorithmic redistricting.}
Minimum-edge-cut redistricting \citep{karypis1998, b1paper} and its extensions
--- edge-weighted bisection (B.2), GeoSection (B.8), AreaSection (B.9) ---
are the foundation of this work.
B.9 establishes the key empirical fact exploited here: the proportionality
gap in competitive states is $\sim$$-6.2$\,pp under pure geographic algorithms,
and this gap is stable across different geographic constraints.
B.11 (PrimeFactor) provides the Huntington-Hill seat allocation $k_D$, $k_R$
used to parameterise the $\mathtt{tpwgts}$ targets.

\paragraph{Multi-constraint graph partitioning.}
METIS $\mathtt{ncon} > 1$ has been used in scientific computing for load
balancing with multiple objectives \citep{karypis1999metis}.
AreaSection (B.9) introduced $\mathtt{ncon}=2$ for redistricting with vertex
weights $[\text{population},\, \text{land area}]$.
This paper applies the same $\mathtt{ncon}=2$ infrastructure with vertex
weights $[\text{population},\, D_{\text{votes}}]$, establishing the partisan
analogue of the area-balance constraint.

\paragraph{Partisan Lorenz curves.}
\citet{rodden2019} uses density plots of partisan composition by geography to
visualise sorting.
\citet{frank2010} applies spatial Lorenz curves to income concentration.
We introduce the \emph{partisan Lorenz curve} --- tracts sorted by Democratic
density, mapping cumulative population to cumulative Democratic-voter fraction ---
as a feasibility diagnostic for proportional packing, directly paralleling the
population-area Lorenz curve of B.9.

\paragraph{Legal background.}
\emph{Rucho v.\ Common Cause} (2019) held that federal courts cannot remedy
partisan gerrymandering under the federal Constitution.
State courts remain empowered to do so under state constitutions
\citep{lwv2018pa, harper2022nc}.
This paper's analysis of the minimum partisan signal required for proportionality
is directly relevant to state-court proportionality claims: it measures the
\emph{cost} of proportionality in partisan information, which is one factor
courts weigh when evaluating whether geographic sorting constitutes a
constitutional violation.
